\begin{proofbox}

	\textbf{\textcolor{CProof}{思路提示}}

	利用和角公式 $\sin(\alpha+\beta)$ 与 $\cos(\alpha+\beta)$，令 $\alpha=\beta=x$ 直接得到正弦与余弦的二倍角公式；正切二倍角公式则在相应定义域内由商数关系推得。

	\medskip
	\textbf{\textcolor{CProof}{正式推导}}
	\begin{description}[style=nextline, leftmargin=2.8em, labelsep=0.5em, itemsep=0.55em]
		\item[\textbf{步骤一（正弦二倍角）：}]
		      在正弦和角公式
		      \[
			      \sin(\alpha+\beta) = \sin\alpha\cos\beta + \cos\alpha\sin\beta
		      \]
		      中取 $\alpha = \beta = x$，得
		      \[
			      \begin{aligned}
				      \sin 2x & = \sin x \cos x + \cos x \sin x \\
				              & = 2\sin x \cos x.
			      \end{aligned}
		      \]
		      \par\textbf{理由：} 和角公式是已知基本公式，代入等角即得。

		\item[\textbf{步骤二（余弦基本形式）：}]
		      在余弦和角公式
		      \[
			      \cos(\alpha+\beta) = \cos\alpha\cos\beta - \sin\alpha\sin\beta
		      \]
		      中取 $\alpha = \beta = x$，得
		      \[
			      \cos 2x = \cos^2 x - \sin^2 x.
		      \]
		      \par\textbf{理由：} 和角公式直接代入。

		\item[\textbf{步骤三（余弦变形）：}]
		      利用同角关系 $\sin^2 x + \cos^2 x = 1$，将 $\cos^2 x$ 或 $\sin^2 x$ 代换，得到另外两种形式：
		      \[
			      \cos 2x = 1 - 2\sin^2 x.
		      \]
		      \[
			      \cos 2x = 2\cos^2 x - 1.
		      \]
		      \par\textbf{理由：} 由 $\cos^2 x = 1 - \sin^2 x$ 或 $\sin^2 x = 1 - \cos^2 x$ 代入基本形式。

		\item[\textbf{步骤四（正切二倍角）：}]
		      当
		      \[
			      \cos 2x\ne0,\qquad \cos x\ne0
		      \]
		      时，由商数关系可得
		      \[
			      \begin{aligned}
				      \tan 2x
				       & = \frac{\sin 2x}{\cos 2x}                  \\
				       & = \frac{2\sin x\cos x}{\cos^2 x-\sin^2 x}.
			      \end{aligned}
		      \]
		      分子分母同除以 $\cos^2 x$，得
		      \[
			      \tan 2x = \frac{2\tan x}{1-\tan^2 x}.
		      \]
		      其中
		      \[
			      1-\tan^2x\ne0
		      \]
		      等价于 $\cos 2x\ne0$，保证分母不为零。
		      \par\textbf{理由：} 商数关系要求 $\cos 2x\ne0$，弦化切要求 $\cos x\ne0$。
	\end{description}

	\medskip
	\textbf{\textcolor{CProof}{结论回扣}}

	\textbf{正弦、余弦二倍角公式对任意实数 $x$ 成立；正切二倍角公式是在 $\cos x\ne0$ 且 $1-\tan^2x\ne0$ 的条件下成立。}

\end{proofbox}